\documentclass[a4paper,12pt]{article}
\newcounter{CJKLanguage}
\setcounter{CJKLanguage}{3}
\input{LaTeX-ToolKit/template.tex}
\begin{document}
\title{日本語文法}
\author{沈威宇}
\date{\temtoday}
\titletocdoc
\sct{日本語文法}
\ssc{コソアド}
\begin{longtable}[c]{|C{0.16\textwidth}|C{0.16\textwidth}|C{0.16\textwidth}|C{0.16\textwidth}|C{0.16\textwidth}|}
\hline
丁寧/普通 & 話し手の近く & 聞き手の近く & それ以外 & 疑問語 \\\hline
連体詞名詞修飾 & この & その & あの & どの \\\hline
連体詞属性 & こんな & そんな & あんな & どんな \\\hline
指示代名詞事物 & こちら/これ & そちら/それ & あちら/あれ & どちら（also, 3つ以上の選択）/どれ（also, 2つの選択） \\\hline
指示代名詞人 & こちらの方、こちら、この方/この人、こいつ & そちらの方、そちら、その方/その人、そいつ & あちらの方、あちら、あの方/あの人、あいつ & どちらの方、どちら、どの方/どの人、どいつ \\\hline
指示代名詞場所 & こちら/ここ & そちら/そこ & あちら/あそこ & どちら/どこ \\\hline
指示代名詞方向 & こちら/こっち & そちら/そっち & あちら/あっち & どちら/どっち \\\hline
副詞 & こう & そう & あう & どう \\\hline
\end{longtable}
\ssc{過去形}
\begin{longtable}[c]{|C{0.15\textwidth}|C{0.13\textwidth}|C{0.13\textwidth}|C{0.13\textwidth}|C{0.13\textwidth}|C{0.13\textwidth}|}
\hline
 & 一段（いちだん）活用（かつよ）動詞（どうし） & イ形容詞（けいようし）/形容詞 & イ形容詞\ruby{良}{い}い & ナ形容詞/形容動詞（けいようどうし） & 名詞（めいし） \\\hline
普通形非過去形肯定 & 〜る & 〜い & \ruby{良}{い}い & 〜だ & 〜だ \\\hline
丁寧形非過去形肯定 & 〜ます & 〜いです & \ruby{良}{い}いです & 〜です & 〜です \\\hline
普通形非過去形否定 & 〜ない & 〜くない & \ruby{良}{よ}くない & 〜じゃない/〜ではない & 〜じゃない/〜ではない \\\hline
丁寧形非過去否定 & 〜ません & 〜くないです/〜くありません & \ruby{良}{よ}くないです/\ruby{良}{よ}くありません & 〜じゃありません/〜ではありません & 〜じゃありません/〜ではありません \\\hline
普通形過去形肯定 & 〜た & 〜かった & \ruby{良}{よ}かった & 〜だった & 〜だった \\\hline
丁寧形過去形肯定 & 〜ました & 〜かったです & \ruby{良}{よ}かったです & 〜でした & 〜でした \\\hline
普通形過去形否定 & 〜なかった & 〜くなかった & \ruby{良}{よ}くなかった & 〜じゃなかった/〜ではなかった & 〜じゃなかった/〜ではなかった \\\hline
丁寧形過去形否定 & 〜ませんでした & 〜くなかったです/〜くありませんでした & \ruby{良}{よ}くなかったです/\ruby{良}{よ}くありませんでした & 〜じゃなかったです/〜じゃありませんでした/〜ではなかったです/〜ではありませんでした & 〜じゃなかったです/〜じゃありませんでした/〜ではなかったです/〜ではありませんでした \\\hline
\end{longtable}
\ssc{助詞}
\begin{longtable}[c]{|C{0.1\textwidth}|C{0.1\textwidth}|C{0.2\textwidth}|C{0.2\textwidth}|C{0.2\textwidth}|}
\hline
助詞 & 種類 & 働き & 文型 & 例文 \\\hline
は & 副助詞 & 主題 & 主題名詞 は 解說 & 私は学生です。 \\\hline
は & 副助詞 & 対比 & 対比名詞 は 解說 & コーヒーは飲みますが、紅茶\textbf{は}飲みません。 \\\hline
も & 副助詞 & 添加 & 添加名詞 は 解說 & 私も学生です。 \\\hline
が & 格助詞 & 主語 & 主語名詞 が 動詞 & 友達が来ました。 \\\hline
が & 格助詞 & 対象（自動詞） & 対象名詞 が 自動詞 & 寿司が食べます。 \\\hline
を & 格助詞 & 対象（他動詞） & 対象名詞 を 他動詞 & ジュースを買います。 \\\hline
に & 格助詞 & 時間 & 時間 に 動詞 & 8時に起きる。 \\\hline
か & 終助詞 & 疑問 & 〜か & 今、何時ですか。 \\\hline
か & 終助詞 & 反語 & 〜か & なぜあなたを忘れようか。 \\\hline
か/よ/ね & 終助詞 & 感動 & 〜か/よ/ね & 合格したか/よ/ね。おめでとう。 \\\hline
ね & 終助詞 & 確認を求める & 〜ね & 今日は暑いですね。 \\\hline
ね & 終助詞 & 答えとして確認を表す & 〜ね & 今日は暑いですね。そうだ\textbf{ね}。 \\\hline
よね & 終助詞 & 確認を求める（確信がもてない） & 〜よね & この映画、おもしろかったよね。 \\\hline
よ/ね & 終助詞 & 告知 & 〜よ/ね & 10時からですよ/ね。 \\\hline
よ & 終助詞 & 勧誘 & 〜よ & いっしょに行こうよ。 \\\hline
よ & 終助詞 & 命令 & 〜よ & 危ないよ。 \\\hline
と & 並列助詞 & 全て挙げる & N1 と N2 と … Nn & コーヒーと紅茶と牛乳がある。 \\\hline
も & 並列助詞 & 全て挙げる & N1 も N2 も … Nn も & コーヒーも紅茶も牛乳もがある。 \\\hline
か & 並列助詞 & 選択（１つを選ぶ） & N1 か N2 & コーヒーか紅茶か牛乳がある。 \\\hline
や & 並列助詞 & 全て挙げない & N1 や N2 や … Nn（や can be replaced with 、 except the last one） & コーヒーや紅茶や牛乳がある。/コーヒー、紅茶や牛乳がある。 \\\hline
とか & 並列助詞 & 全て挙げない（二つ目は必要ではない） & N1 とか N2 とか … Nn とか & コーヒーとか紅茶とか牛乳とかがある。 \\\hline
だの & 並列助詞 & 全て挙げない（二つ目は必要ではない、話し手の不満の感情が含まれる） & N1 だの N2 だの … Nn だの & コーヒーだの紅茶だの牛乳だのがある。 \\\hline
\end{longtable}
\ssc{数詞}
Attachments of 一:
\begin{itemize}
\item Positive integer powers of 万 begin with 一 if no digit would otherwise precede.
\item If 千 directly precedes the name of positive integer powers of 万, 一 is attached before 千, which yields \ruby{一千}{いっせん}.
\item If 千 does not directly precede the name of positive integer powers of 万, attaching 一 is optional.
\item 一 is normally not attached before 百 or 十.
\end{itemize}
Phonetic modifications and reading conventions:
\begin{itemize}
\item 四 is typically pronounced よん instead of し; 七 is typically pronounced なな instead of しち.
\item \ruby{何}{なん}'s phoentic modifications follow \ruby{三}{さん}'s.
\item \ruby{百}{ひゃく}: いっぴゃく、さんびゃく、ろっぴゃく、はっぴゃく。
\item \ruby{千}{せん}: いっせん、さんぜん、はっせん。
\item \ruby{万}{まん}: None.
\item \ruby{億}{おく}: None.
\item \ruby{兆}{ちょう}: いっちょう、はっちょう、じゅっちょう。
\item \ruby{時}{じ}: よじ、しちじ、くじ。
\item \ruby{分}{ふん}: いっぷん、さんぷん、よんぷん、ろっぷん、はっぷん、じゅっぷん。
\item \ruby{秒}{びょう}: None.
\end{itemize}
Strings of digits:
\begin{itemize}
\item Lengthen normally unimoraic digits such as に, ご, and less used, し and く to form groups of two digits of exactly four morae as long as possible, e.g., 五四二六二 is pronounced ごうよんにいろくに, and 二点二五五 is pronounced にいてんにいごうご.
\item - between digits is pronounced の.
\item 零 is also read as まる and written as 〇.
\end{itemize}
\ssc{時間を表す言葉}
\begin{itemize}
\item \ruby{一昨々日}{さきおととい} $\rightarrow$ \ruby{一昨日}{おととい} $\rightarrow$ \ruby{昨日}{きのう} $\rightarrow$ \ruby{今日}{きょう} $\rightarrow$ \ruby{明日}{あした} $\rightarrow$ \ruby{明後日}{あさって} $\rightarrow$ \ruby{明々後日}{しあさって}
\item \ruby{先々週}{せんせんしゅう} $\rightarrow$ \ruby{先週}{せんしゅう} $\rightarrow$ \ruby{今週}{こんしゅう} $\rightarrow$ \ruby{来週}{らいしゅう} $\rightarrow$ \ruby{再来週}{さらいしゅう}
\item \ruby{朝}{あさ}: \ruby{今朝}{けさ}。
\item \ruby{昼}{ひる}: None.
\item \ruby{夕方}{ゆうがた}: None.
\item \ruby{夜}{よる}/\ruby{晩}{ばん}: \ruby{昨夜}{ゆうべ}、\ruby{今晩}{こんばん}/\ruby{今夜}{こんや}。
\end{itemize}

\ssc{動詞の活用}
TODO

五段活用/グループ1

（上/下）一段活用/グループ2

来る

する
\end{document}
